\chapter{Dostateczność, zupełność i optymalna estymacja nieobciążona}
\label{ch:sufficiency}
\paragraph{Powiązania i odsyłacze.} Punktem wyjścia są statystyki i estymatory z rozdziałów~\ref{ch:statystyki} i~\ref{ch:estimators}; ograniczenia efektywności opisuje rozdział~\ref{ch:fisher-information}.


\section{Czy wolno zastąpić całą próbę jedną liczbą?}
Dla $n$ niezależnych prób Bernoulliego zapisujemy wektor
$X=(X_1,\ldots,X_n)$, ale estymujemy parametr jedynie przez
liczbę sukcesów $S=\sum_i X_i$. Dlaczego wolno wyrzucić
informację o kolejności wyników? W tym konkretnym modelu,
\emph{warunkowo na $S=k$}, każda sekwencja zawierająca $k$
sukcesów ma to samo prawdopodobieństwo $\binom nk^{-1}$,
które \emph{nie zależy od $p$}. Znajomość kolejności nie wnosi
więc nowej informacji o parametrze. Taką redukcję formalizuje
dostateczność. Nie każda statystyka ma tę własność: w rozkładzie
jednostajnym ze zmiennym kresem informację skupia maksimum,
a w ogólnym modelu zależnym od czasu kolejność może być istotna.

\begin{definicja}[Statystyka dostateczna]\label{def:sufficiency}
Niech $X$ ma prawo $P_\theta$ na przestrzeni standardowej borelowskiej
$(\mathsf X,\calA)$ dla każdego $\theta\in\Theta$,
a $T:\mathsf X\to\mathsf T$ będzie mierzalna, przy czym
$\mathsf T$ także jest przestrzenią standardową borelowską.
Mówimy, że $T$ jest \emph{dostateczna} dla rodziny
$\{P_\theta\}$, jeśli istnieje jądro probabilistyczne
$K(t,A)$, mierzalne w $t$ dla każdego $A\in\calA$,
takie, że dla wszystkich $\theta$ i $A\in\calA$
\[
P_\theta(X\in A\mid\sigma(T))=K(T,A)
\quad P_\theta\text{-prawie na pewno}.
\]
\end{definicja}

Warunkowy rozkład pełnych danych przy ustalonej statystyce musi mieć
\emph{jedną wspólną wersję} niezależną od parametru. Nie wystarcza,
że dla każdego parametru oddzielnie istnieje pewien rozkład
warunkowy: to zawsze jest prawdą w przestrzeniach standardowych
borelowskich. Dostateczność jest własnością \emph{rodziny rozkładów},
a nie jednej obserwacji.

\begin{zadanie}[Dlaczego liczba sukcesów wystarcza?]\label{zad:suff-bernoulli}
Niech $X_1,\ldots,X_n\iid\operatorname{Bernoulli}(p)$,
$p\in[0,1]$, $S=\sum_iX_i$. Oblicz prawo pełnej sekwencji
warunkowo na $S=k$.
\par\emph{Rozwiązanie.} Dla $0<p<1$ i sekwencji
$x\in\{0,1\}^n$ o sumie $k$ wzór Bayesa daje
\[
P_p(X=x\mid S=k)=\frac{p^k(1-p)^{n-k}}
 {\binom nk p^k(1-p)^{n-k}}=\binom nk^{-1}.
\]
Dla sekwencji o innej sumie prawdopodobieństwo jest zerowe.
Przy $p=0$ obserwujemy $S=0$ prawie na pewno, a przy $p=1$
$S=n$ prawie na pewno; ta sama formuła zdefiniowana dla wszystkich
$k$ nadal stanowi wspólną wersję jądra. Jest więc ono niezależne
od $p$, zatem $S$ jest dostateczna. Wniosek praktyczny:
w modelu \emph{iid Bernoulliego} do wnioskowania o $p$ można
zachować tylko $S$ i $n$; jeśli prawdopodobieństwo sukcesu zmienia
się z czasem, redukcja może zgubić informację o trendzie.
\end{zadanie}


\begin{center}
\begin{tikzpicture}[>=Latex,node distance=10mm,
 box/.style={draw,rounded corners,align=center,inner sep=5pt,
 font=\small,minimum height=13mm}]
\node[box,text width=30mm] (full) {Pełna próba $X$\\kolejność i wyniki};
\node[box,text width=38mm,right=of full] (stat) {Statystyka $T(X)$\\informacja o parametrze};
\node[box,text width=41mm,right=of stat] (cond) {Warunkowe prawo $X\mid T$\\nie zależy od $\theta$};
\draw[->,thick] (full)--(stat);
\draw[->,thick] (stat)--(cond);
\end{tikzpicture}
\end{center}
\noindent\emph{Rysunek 6.1.} Dostateczność nie oznacza możliwości
odtworzenia pełnej próby z $T$. Oznacza, że pozostała losowość
próby po ustaleniu $T$ nie zależy od nieznanego parametru.

\section{Dlaczego istnieje wspólna miara odniesienia?}
Aby uzasadnić kryterium faktoryzacyjne bez pomijania warunków
mierzalności, przypomnimy jeden lemat. W całym podrozdziale
zakładamy standardowe przestrzenie borelowskie i zdominowaną
rodzinę $P_\theta\ll\mu$, gdzie $\mu$ jest $\sigma$-skończona.
Jądra warunkowe względem $T$ istnieją w tych przestrzeniach,
co przyjmujemy jako standardowy wynik o regularnych rozkładach
warunkowych.

\begin{lemat}[Przeliczalna mieszanina dominująca]\label{lem:halmos-savage}
Dla dowolnej rodziny miar probabilistycznych $\{P_\theta\}_{\theta\in\Theta}$
zdominowanej przez jedną miarę $\sigma$-skończoną $\mu$
istnieją parametry $\theta_1,\theta_2,\ldots$ oraz dodatnie wagi
$a_j$ o sumie $1$ takie, że prawdopodobieństwo
$P_*:=\sum_{j\ge1}a_jP_{\theta_j}$ dominuje każde $P_\theta$.
\end{lemat}
\begin{proof}
Wybierzmy miarę skończoną $\nu$ równoważną $\mu$: jeśli
$\mathsf X=\bigcup_{j\ge1}A_j$ i $0<\mu(A_j)<\infty$
(po usunięciu zbiorów zerowych i rozłączeniu), możemy położyć
$\nu(B)=\sum_j2^{-j}\mu(B\cap A_j)/(1+\mu(A_j))$,
a następnie znormalizować. Dla każdego $\theta$ niech
$B_\theta=\{\dd P_\theta/\dd\mu>0\}$.
Niech $c$ będzie supremum $\nu$-miar skończonych sum
$\bigcup_{i=1}^m B_{\theta_i}$. Dla każdego $k$ wybieramy
skończoną rodzinę, której suma $D_k$ spełnia
$\nu(D_k)>c-1/k$. Po zebraniu wszystkich wybranych parametrów
w jeden ciąg otrzymujemy $D=\bigcup_kD_k$ z $\nu(D)=c$
(z jednej strony $D$ jest rosnącą granicą skończonych sum,
każda o mierze nie większej niż $c$). Dla każdego $\theta$,
gdyby $\nu(B_\theta\setminus D)>0$, odpowiednio duża skończona
część $D$ powiększona o $B_\theta$ miałaby miarę większą niż $c$,
sprzeczność. Stąd $\nu(B_\theta\setminus D)=0$.
Niech $P_*=\sum_j a_jP_{\theta_j}$, np. $a_j=2^{-j}$.
Jego gęstość jest dodatnia na $D$ z dokładnością do zbioru
$\mu$-zerowego. Jeżeli $P_*(A)=0$, to
$\mu(A\cap D)=0$, a zatem dla każdego $\theta$
$P_\theta(A)=P_\theta(A\cap D)=0$. To oznacza $P_\theta\ll P_*$.
\end{proof}

\begin{twierdzenie}[Kryterium faktoryzacyjne Fishera--Neymana]
\label{thm:factorization}
Niech $\{P_\theta:\theta\in\Theta\}$ będzie rodziną miar na
standardowej przestrzeni borelowskiej $(\mathsf X,\calA)$,
zdominowaną przez miarę $\sigma$-skończoną $\mu$.
Niech $p_\theta=\dd P_\theta/\dd\mu$ będą wybranymi wersjami
gęstości. Statystyka $T:\mathsf X\to\mathsf T$ o wartościach
w przestrzeni standardowej borelowskiej jest dostateczna wtedy
i tylko wtedy, gdy istnieją mierzalne funkcje $h\ge0$ na
$\mathsf X$ i $g_\theta\ge0$ na $\mathsf T$ takie, że
\[
p_\theta(x)=g_\theta(T(x))h(x)
\quad\mu\text{-prawie wszędzie dla każdego }\theta.
\]
Zbiór, na którym równość zawodzi, może zależeć od $\theta$.
\end{twierdzenie}
\begin{proof}
Niech $P_*$ będzie mieszaniną z lematu~\ref{lem:halmos-savage},
a $p_*$ jej gęstością względem $\mu$. Dla każdego $\theta$
mamy $P_\theta\ll P_*$ i pochodną
$r_\theta=\dd P_\theta/\dd P_*$. Załóżmy najpierw,
że $r_\theta(x)=g_\theta(T(x))$ $P_*$-prawie wszędzie.
Niech $K(t,A)$ będzie wspólną wersją regularnego prawa warunkowego
$X$ względem $T$ pod $P_*$. Dla każdego $A\in\calA$
i mierzalnego $B\subset\mathsf T$, korzystając z zależności
$r_\theta=g_\theta(T)$, otrzymujemy
\[
P_\theta(A\cap\{T\in B\})
=\E_*[\ind{\{X\in A\}}g_\theta(T)\ind{\{T\in B\}}]
=\E_*[K(T,A)g_\theta(T)\ind{\{T\in B\}}]
=\E_\theta[K(T,A)\ind{\{T\in B\}}].
\]
Wobec definicji warunkowania $K$ jest wspólnym jądrem warunkowym,
więc $T$ jest dostateczna.

Odwrotnie, jeśli $T$ jest dostateczna z jądrem $K$, dla każdej
ograniczonej funkcji mierzalnej $f$ mamy
$\E_\theta f(X)=\E_\theta\int_{\mathsf X}f(x)K(T,\dd x)$.
Ta sama równość zachodzi pod mieszaniną $P_*$ (sumujemy ją
z wagami $a_j$); oznacza to, że
$\int_{\mathsf X}f(x)K(T,\dd x)=\E_*[f(X)\mid T]$ $P_*$-p.n.
Ponieważ $P_\theta\ll P_*$, dla każdego ograniczonego $f$
\[
\E_*[f(X)r_\theta(X)]
=\E_\theta f(X)
=\E_\theta\E_*[f(X)\mid T]
=\E_*[r_\theta(X)\E_*[f(X)\mid T]].
\]
Z własności warunkowej wartości oczekiwanej ostatnie wyrażenie
równa się $\E_*[f(X)\E_*(r_\theta(X)\mid T)]$.
Dowolność ograniczonego $f$ implikuje
$r_\theta=\E_*(r_\theta\mid T)$ $P_*$-p.n., więc z lematu
Dooba--Dynkina $r_\theta=g_\theta(T)$ $P_*$-p.n.
Równoważność z podanym rozkładem gęstości wynika następująco.
Z zależności $p_\theta=r_\theta p_*$ dostajemy faktoryzację
z $h=p_*$. Jeżeli zaś
$p_\theta=g_\theta(T)h$, to mieszanka ma gęstość
$p_*=(\sum_j a_jg_{\theta_j}(T))h$.
Na zbiorze $p_*>0$ iloraz $p_\theta/p_*$ zależy tylko od $T$,
więc jest $\sigma(T)$-mierzalny. Na zbiorze zerowym $P_*$
wersję ilorazu definiujemy dowolnie. Oba kierunki są dowiedzione.
\end{proof}

\begin{przyklad}[Trzy modele i trzy różne statystyki dostateczne]
\begin{enumerate}
\item W próbie Bernoulliego statystyką dostateczną jest
$S=\sum_iX_i$, ponieważ dla obserwacji $x$ mamy
\[
p_p(x)=p^{S(x)}(1-p)^{n-S(x)}.
\]
\item Poisson: dla $X_i\iid\Pois(\lambda)$,
$p_\lambda(x)=\exp(-n\lambda)\lambda^{\sum_i x_i}
\prod_i(x_i!)^{-1}$; wystarczy $S=\sum_iX_i$.
\item Normalny: dla próby $\Normal(\mu,\sigma^2)$ z \emph{dwoma}
nieznanymi parametrami gęstość zależy od danych przez
$S_1=\sum_iX_i$ i $S_2=\sum_iX_i^2$.
Kiedy $\sigma^2$ jest znane, wystarcza samo $\bar X$.
\end{enumerate}
\end{przyklad}

\section{Minimalna dostateczność: kiedy nie zachowaliśmy za dużo?}
Dostateczną statystykę można zawsze sztucznie rozszerzyć:
wektor $(S,X_1)$ jest dostateczny, jeżeli $S$ jest dostateczny.
Chcielibyśmy zachować \emph{jak najmniej} informacji niezbędnej
do opisu parametru. W modelach z dodatnimi gęstościami na wspólnym
nośniku prowadzi do tego proste kryterium ilorazów wiarygodności.

\begin{definicja}[Minimalna dostateczność]
Dostateczna statystyka $T$ jest \emph{minimalnie dostateczna},
jeśli dla każdej innej dostatecznej statystyki $S$ istnieje mierzalna
funkcja $\varphi$ taka, że $T=\varphi(S)$ $P_\theta$-prawie na pewno
dla każdego $\theta$.
\end{definicja}

\begin{stwierdzenie}[Kryterium ilorazowe w modelu dyskretnym]
\label{prop:minimal-discrete}
Niech $\mathsf X$ będzie przeliczalna i niech wszystkie funkcje
masy $p_\theta(x)$ będą \emph{ściśle dodatnie} dla $x\in\mathsf X$.
Załóżmy, że istnieje statystyka $T$ spełniająca
\[
T(x)=T(y)\quad\Longleftrightarrow\quad
\frac{p_\theta(x)}{p_\theta(y)}\text{ nie zależy od }\theta.
\]
Wówczas $T$ jest minimalnie dostateczna.
\end{stwierdzenie}
\begin{proof}
Wybierzmy ustalone $\theta_0$. Jeśli $T(x)=T(y)$, to
$p_\theta(x)/p_\theta(y)=p_{\theta_0}(x)/p_{\theta_0}(y)$.
W każdej klasie wartości $t$ wybierzmy reprezentanta $x_t$
(co jest możliwe dla przeliczalnej przestrzeni). Definiując
$g_\theta(t)=p_\theta(x_t)/p_{\theta_0}(x_t)$ oraz
$h(x)=p_{\theta_0}(x)$ dostajemy
$p_\theta(x)=g_\theta(T(x))h(x)$, zatem $T$ jest dostateczna
z twierdzenia~\ref{thm:factorization}.
Niech teraz $S$ będzie dowolną inną dostateczną statystyką.
Faktoryzacja daje
$p_\theta(x)=u_\theta(S(x))v(x)$, gdzie $v(x)>0$
na całej przeliczalnej przestrzeni (po ewentualnej zmianie
wersji na zbiorach zerowych: takich punktów nie ma, bo wszystkie
$p_\theta$ są dodatnie). Zatem $S(x)=S(y)$ implikuje
$p_\theta(x)/p_\theta(y)=v(x)/v(y)$, niezależnie od $\theta$,
a więc $T(x)=T(y)$. To znaczy, że $T$ jest stała na włóknach
$S$; dla przeliczalnej dziedziny definiujemy mierzalną funkcję
$\varphi$ na zbiorze wartości $S$ przez $\varphi(S(x))=T(x)$,
a poza nim dowolnie. Otrzymujemy $T=\varphi(S)$.
\end{proof}

\begin{przyklad}[W modelu Bernoulliego liczba sukcesów jest minimalnie dostateczna]
Niech $p\in(0,1)$, $x,y\in\{0,1\}^n$. Wówczas
\[
\frac{p_p(x)}{p_p(y)}
=\left(\frac p{1-p}\right)^{\sum_i x_i-\sum_i y_i}.
\]
Iloraz jest niezależny od $p$ dokładnie wtedy, gdy
$\sum_i x_i=\sum_i y_i$. Zatem liczba sukcesów $S$ jest
minimalnie dostateczna. Oznacza to minimalność w sensie
relacji mierzalności, \emph{nie} najmniejszej liczby bitów w pliku
ani automatycznie minimalnego błędu estymacji.
\end{przyklad}

\section{Rao--Blackwell: wykorzystanie informacji bez dodawania danych}
Wyobraźmy sobie, że estymujemy $p$ na podstawie pierwszego wyniku
$X_1$, choć mamy do dyspozycji całą próbę. Statystyka dostateczna
$S$ pozwala zastąpić ten niestabilny estymator jego \emph{średnią
warunkową przy znanym $S$}. Prowadzi to do mniejszej wariancji:
operacja warunkowania usuwa przypadkowość, która przy znanym $S$
nie niesie informacji o parametrze.

\begin{twierdzenie}[Twierdzenie Rao--Blackwella]\label{thm:rao-blackwell}
Niech $T$ będzie dostateczną statystyką dla rodziny
$\{P_\theta:\theta\in\Theta\}$ na standardowych przestrzeniach
borelowskich, a $U(X)$ rzeczywistym estymatorem takim, że
$\E_\theta|U|<\infty$ dla każdego $\theta$. Istnieje statystyka
$U^*=u(T)$, \emph{niezależna od wyboru parametru}, taka że
$U^*=\E_\theta[U\mid T]$ $P_\theta$-p.n. dla każdego $\theta$.
Jeśli $\E_\theta U=g(\theta)$, to także
$\E_\theta U^*=g(\theta)$.
Dla każdej wypukłej funkcji straty $a\mapsto L(\theta,a)$,
przy której obie wartości oczekiwane są określone,
\[
\E_\theta L(\theta,U^*)\le\E_\theta L(\theta,U).
\]
Dla straty kwadratowej, jeśli $U\in L^2(P_\theta)$, zachodzi
ściślejsza tożsamość
\[
\E_\theta(U-g(\theta))^2
=\E_\theta(U^*-g(\theta))^2
 +\E_\theta\bigl(U-U^*\bigr)^2.
\]
\end{twierdzenie}
\begin{proof}
Z dostateczności istnieje wspólne jądro warunkowe $K(t,\dd x)$.
Definiujemy $u(t)=\int_{\mathsf X}U(x)K(t,\dd x)$ w punktach, gdzie całka
jest bezwzględnie zbieżna, a poza nimi przyjmujemy $u(t)=0$.
Funkcja $u$ jest mierzalna (najpierw dla funkcji prostych,
potem przez aproksymację i zbieżność monotoniczną).
Ponieważ $U\in L^1(P_\theta)$ dla każdego $\theta$,
$u(T)=\E_\theta(U\mid T)$ prawie na pewno pod każdym $P_\theta$.
Z własności wieży $\E_\theta U^*=\E_\theta U$.
Warunkowa nierówność Jensena daje
$L(\theta,\E_\theta(U\mid T))\le
\E_\theta(L(\theta,U)\mid T)$;
po scałkowaniu otrzymujemy nierówność ryzyk.
Dla straty kwadratowej piszemy
$U-g(\theta)=(U-U^*)+(U^*-g(\theta))$.
Człon mieszany ma wartość oczekiwaną zero, bo
$\E_\theta[U-U^*\mid T]=0$ i $U^*-g(\theta)$ jest
$\sigma(T)$-mierzalny oraz całkowalny w $L^2$.
Rozwinięcie kwadratu daje tezę.
\end{proof}

\begin{zadanie}[Od pierwszej obserwacji do udziału sukcesów]
Dla $X_i\iid\operatorname{Bernoulli}(p)$ zacznij od
$U=X_1$, który estymuje $p$ nieobciążenie.
Wyznacz $\E(U\mid S)$, gdzie $S=\sum_iX_i$, i porównaj wariancje.
\par\emph{Rozwiązanie.} Warunkowo na $S=k$ każda pozycja
ma prawdopodobieństwo sukcesu $k/n$ z symetrii sekwencji,
więc $\E(X_1\mid S)=S/n=\bar X$.
Zatem $\Var_p(U)=p(1-p)$,
a $\Var_p(U^*)=p(1-p)/n$.
Zysk jest szczególnie duży, gdy $n$ rośnie: nie dodaliśmy
żadnej nowej obserwacji, wykorzystaliśmy tylko \emph{informację
już obecną w całej próbie}.
\end{zadanie}

\section{Zupełność: dlaczego nie ma dwóch różnych UMVU?}
Dostateczność mówi, że $T$ zachowuje informację o parametrze.
Zupełność zapewnia coś innego: nie istnieje niezerowa funkcja
$T$ o wartości oczekiwanej identycznie równej zero dla wszystkich
parametrów. To właśnie gwarantuje jednoznaczność estymatora
nieobciążonego opartego na statystyce dostatecznej.

\begin{definicja}[Zupełność rodziny statystyki]\label{def:completeness}
Rodzina rozkładów $\{\Law_\theta(T):\theta\in\Theta\}$
jest \emph{zupełna}, jeżeli dla każdej mierzalnej $a$ takiej, że
$\E_\theta|a(T)|<\infty$ dla wszystkich $\theta$, z warunku
$\E_\theta a(T)=0$ dla wszystkich $\theta$ wynika
$a(T)=0$ $P_\theta$-prawie na pewno dla wszystkich $\theta$.
\end{definicja}

\begin{twierdzenie}[Zupełność liczby sukcesów Bernoulliego]
\label{thm:complete-binomial}
Niech $S=\sum_{i=1}^nX_i$, gdzie
$X_i\iid\operatorname{Bernoulli}(p)$, $p\in(0,1)$.
Wtedy rodzina rozkładów $S$ jest zupełna.
\end{twierdzenie}
\begin{proof}
Niech $a:\{0,\ldots,n\}\to\R$ spełnia
$\E_p a(S)=0$ dla każdego $p\in(0,1)$. Mamy
\[
0=\sum_{k=0}^n a(k)\binom nk p^k(1-p)^{n-k}.
\]
Dzieląc przez $(1-p)^n>0$ i podstawiając
$z=p/(1-p)\in(0,\infty)$, otrzymujemy
$\sum_{k=0}^n a(k)\binom nk z^k=0$ dla każdego $z>0$.
Wielomian mający nieskończenie wiele zer jest zerowy,
więc wszystkie jego współczynniki znikają:
$a(k)\binom nk=0$. Ponieważ $\binom nk>0$,
$a(k)=0$ dla wszystkich $k$. Dowiedliśmy nawet silniejszego
wniosku niż wymagana równość prawie na pewno.
\end{proof}

\begin{twierdzenie}[Zupełność sumy zmiennych Poissona]
\label{thm:complete-poisson}
Niech $X_i\iid\Pois(\lambda)$, $\lambda\in(0,\infty)$,
a $S=\sum_iX_i\sim\Pois(n\lambda)$. Wtedy rodzina praw
$S$ jest zupełna.
\end{twierdzenie}
\begin{proof}
Jeśli $\E_\lambda|a(S)|<\infty$ i
$\E_\lambda a(S)=0$ dla każdego $\lambda>0$, to
\[
\sum_{k=0}^\infty a(k)\frac{(n\lambda)^k}{k!}=0,
\qquad\lambda>0.
\]
Szereg potęgowy po lewej stronie jest absolutnie zbieżny
dla każdego dodatniego argumentu, ponieważ odpowiednie
wartości bezwzględne są równe
$e^{n\lambda}\E_\lambda|a(S)|<\infty$.
Ma więc nieskończony promień zbieżności i reprezentuje
funkcję analityczną. Skoro znika na $(0,\infty)$,
wszystkie jej współczynniki Taylora w zerze muszą być zerowe.
To daje $a(k)=0$ dla każdego $k$.
\end{proof}


\begin{twierdzenie}[Zupełność pełnej jednoparametrowej rodziny wykładniczej]
\label{thm:complete-exponential-family}
Niech $Q_\eta$, $\eta\in J$, będzie rodziną rozkładów rzeczywistej
statystyki $T$ o postaci
$\dd Q_\eta(t)=e^{\eta t-A(\eta)}\,\dd\rho(t)$,
gdzie $\rho$ jest nieujemną miarą $\sigma$-skończoną,
$J\subset\R$ niepustym \emph{otwartym przedziałem}, a
$0<\int_{\R}e^{\eta t}\,\dd\rho(t)<\infty$ dla każdego $\eta\in J$.
Wtedy rodzina $\{Q_\eta:\eta\in J\}$ jest zupełna.
\end{twierdzenie}
\begin{proof}
Niech $a$ będzie mierzalna, $\E_\eta|a(T)|<\infty$
i $\E_\eta a(T)=0$ dla wszystkich $\eta\in J$.
Ustalmy $\eta_0\in J$ i wybierzmy $d>0$ z
$[\eta_0-2d,\eta_0+2d]\subset J$.
Zdefiniujmy skończoną miarę ze znakiem
$\nu(\dd t)=a(t)Q_{\eta_0}(\dd t)$. Dla $|u|\le d$
całka $\int_{\R}e^{u t}|\nu|(\dd t)$ jest skończona:
na $t\ge0$ ograniczamy ją przez stałą razy całkę
$|a(t)|e^{dt}$ względem $Q_{\eta_0}$,
a na $t<0$ analogicznie przez $e^{-dt}$;
obie są skończone dzięki całkowalności $a$ przy
$\eta_0+d$ i $\eta_0-d$. Funkcja
$F(z)=\int_{\R}e^{zt}\,\nu(\dd t)$ jest więc holomorficzna
w pasku $|\operatorname{Re}z|<d$; różniczkowanie pod całką
uzasadniają analogiczne oszacowania na mniejszych paskach.
Dla rzeczywistego $u\in(-d,d)$ mamy
\[
F(u)=e^{A(\eta_0+u)-A(\eta_0)}\E_{\eta_0+u}a(T)=0.
\]
Z twierdzenia o identyczności funkcji holomorficznych
$F(z)=0$ na całym pasku. W szczególności
$F(iv)=\int_{\R}e^{ivt}\,\nu(\dd t)=0$ dla każdego $v\in\R$.
Z jednoznaczności transformaty Fouriera skończonej miary
ze znakiem wynika $\nu=0$, czyli $a=0$ $Q_{\eta_0}$-p.n.
Gęstości $Q_\eta$ względem $\rho$ są dodatnie na tym samym
nośniku, więc wszystkie $Q_\eta$ są równoważne. Zatem
$a=0$ $Q_\eta$-p.n. dla każdego $\eta\in J$.
\end{proof}
\begin{przyklad}[Zupełność średniej w modelu normalnym]
Gdy $X_i\iid\Normal(\mu,\sigma^2)$ przy znanym $\sigma^2>0$
i $\mu\in\R$, rozkład $T=\sum_iX_i$ jest rodziną wykładniczą
z parametrem naturalnym $\eta=\mu/\sigma^2$ przebiegającym
całą prostą. Z poprzedniego twierdzenia wynika zupełność
statystyki $T$, a z faktoryzacji jej dostateczność.
Wobec tego $\bar X$ jest UMVU dla $\mu$. Nie należy
mieszać tego przypadku z sytuacją, w której również
$\sigma^2$ jest parametrem nieznanym.
\end{przyklad}

\begin{twierdzenie}[Lehmann--Scheffé: jednoznaczność UMVU]
\label{thm:lehmann-scheffe}
Niech $T$ będzie statystyką dostateczną i zupełną dla
$\{P_\theta:\theta\in\Theta\}$, a $g:\Theta\to\R$ estymandem.
Załóżmy, że istnieje estymator $U(X)$ o skończonym drugim
momencie pod każdym $P_\theta$ i taki, że
$\E_\theta U=g(\theta)$ dla każdego $\theta$.
Wtedy $U^*=\E_\theta(U\mid T)$, wybrany jako wspólna
wersja z twierdzenia Rao--Blackwella, jest \emph{jednostajnie
najlepszym nieobciążonym estymatorem} (UMVU): dla każdego
nieobciążonego estymatora $V\in L^2(P_\theta)$
$\Var_\theta(U^*)\le\Var_\theta(V)$ dla wszystkich $\theta$.
Ponadto jest jedyny z dokładnością do równości
$P_\theta$-prawie na pewno dla każdego $\theta$.
\end{twierdzenie}
\begin{proof}
Rao--Blackwell daje $U^*=u(T)$ nieobciążony i o skończonym
momencie kwadratowym. Dla dowolnego innego nieobciążonego
$V$ jego uśrednienie $V^*=v(T)$ jest nieobciążone i
$\Var_\theta V^*\le\Var_\theta V$. Ponieważ
$\E_\theta[u(T)-v(T)]=0$ dla każdego $\theta$ i różnica
jest całkowalna, zupełność implikuje
$u(T)=v(T)$ $P_\theta$-p.n. dla każdego $\theta$.
Stąd $\Var_\theta U^*=\Var_\theta V^*\le\Var_\theta V$.
Jeśli $V$ też jest UMVU, to w szczególności
$\Var_\theta(V)=\Var_\theta(U^*)$ i równość w tożsamości
Pitagorasa Rao--Blackwella wymusza
$V=V^*=U^*$ $P_\theta$-p.n.
\end{proof}

\section{Przykłady UMVU: praktyczne obliczenia}
\begin{zadanie}[Prawdopodobieństwo dwóch sukcesów]
\label{zad:umvu-p2}
Niech $X_1,\ldots,X_n\iid\operatorname{Bernoulli}(p)$,
$n\ge2$. Chcemy oszacować $p^2$ nieobciążenie.
\par\emph{Rozwiązanie.} Estymator $U=X_1X_2$ jest nieobciążony,
gdyż $\E U=\E X_1\E X_2=p^2$.
Niech $S=\sum_iX_i$. Warunkowo na $S=k$ prawdopodobieństwo,
że pierwsze dwie pozycje są sukcesami, wynosi
$k(k-1)/(n(n-1))$: wybieramy dwa z $k$ sukcesów i
porównujemy z wyborem dwóch spośród $n$ pozycji. Wobec tego
\[
\widehat{p^2}_{\rm UMVU}=\E(X_1X_2\mid S)
=\frac{S(S-1)}{n(n-1)}.
\]
Z twierdzeń~\ref{thm:complete-binomial} i~\ref{thm:lehmann-scheffe}
ma najmniejszą wariancję ze wszystkich estymatorów
nieobciążonych $p^2$. Jeżeli $n=4$ i zaobserwowano
$x=(1,0,1,0)$, to $S=2$ i wynik wynosi $1/6$, podczas gdy
wstawienie MLE $\widehat p=1/2$ do $p^2$ daje $1/4$.
To nie sprzeczność: $\widehat p^{\,2}$ jest obciążony jako
estymator $p^2$, ponieważ
$\E\bar X^2=p^2+p(1-p)/n$.
\end{zadanie}

\begin{zadanie}[Prognoza braku zgłoszeń w przyszłym okresie]
\label{zad:umvu-pois-zero}
W $n$ niezależnych okresach obserwujemy liczby zgłoszeń
$X_i\sim\Pois(\lambda)$, $\lambda>0$. Oszacuj nieobciążenie
prawdopodobieństwo braku zgłoszeń w \emph{następnym} takim
okresie, równe $e^{-\lambda}$.
\par\emph{Rozwiązanie.} Statystyka $S=\sum_iX_i\sim\Pois(n\lambda)$
jest dostateczna i zupełna. Funkcja tworząca rozkładu Poissona
daje $\E_\lambda a^S=\exp(n\lambda(a-1))$ dla $a\ge0$.
Dla $a=1-1/n$ (przy $n\ge1$) otrzymujemy
\[
\widehat{e^{-\lambda}}=\left(1-\frac1n\right)^S,
\quad \E_\lambda\widehat{e^{-\lambda}}=e^{-\lambda}.
\]
Dla $n=1$ przyjmujemy $0^0=1$, więc estymator jest wskaźnikiem
$\{S=0\}$. Z twierdzenia Lehmann--Scheffégo jest UMVU.
Przykładowo przy $n=5$ i sumie zgłoszeń $S=3$ wynik to
$(4/5)^3=0{,}512$, natomiast estymator przez wstawienie MLE
$\lambda=S/n$ daje $e^{-3/5}\approx0{,}549$.
Oba są sensowne dla innych kryteriów; UMVU nie gwarantuje
najmniejszego błędu średniokwadratowego wśród \emph{wszystkich}
estymatorów, tylko wśród nieobciążonych.
\end{zadanie}

\section{Statystyki poboczne i twierdzenie Basu}
\begin{definicja}[Statystyka poboczna]
Statystykę $A(X)$ nazywamy \emph{poboczną} (ang. \emph{ancillary}),
jeśli jej rozkład nie zależy od nieznanego parametru $\theta$.
\end{definicja}

\begin{twierdzenie}[Twierdzenie Basu]\label{thm:basu}
Niech $T$ będzie dostateczna i zupełna dla rodziny
$\{P_\theta:\theta\in\Theta\}$ na przestrzeni standardowej
borelowskiej, a $A$ statystyką poboczną o wartościach w przestrzeni
standardowej borelowskiej. Wtedy $T$ i $A$ są niezależne
pod każdym $P_\theta$.
\end{twierdzenie}
\begin{proof}
Ustalmy zdarzenie $B$ z $\sigma$-ciała wartości statystyki $A$.
Z dostateczności istnieje funkcja mierzalna $q_B$ \emph{niezależna
od $\theta$} taka, że
$\E_\theta[\ind{\{A\in B\}}\mid T]=q_B(T)$.
Z poboczności $P_\theta(A\in B)=c_B$ nie zależy od $\theta$.
Zatem dla każdego $\theta$
$\E_\theta[q_B(T)-c_B]=0$.
Funkcja $q_B-c_B$ jest ograniczona, więc zupełność wymusza
$q_B(T)=c_B$ $P_\theta$-p.n. dla wszystkich $\theta$.
Oznacza to, że
$P_\theta(A\in B\mid T)=P_\theta(A\in B)$;
ponieważ zachodzi dla dowolnego $B$, z definicji jest to
niezależność $T$ i $A$.
\end{proof}

\begin{przyklad}[Niezależność średniej i wariancji normalnej]
Dla $X_i\iid\Normal(\mu,\sigma^2)$ z \emph{znanym} $\sigma^2>0$
i nieznaną $\mu\in\R$, statystyka $T=\bar X$ jest
w tym pełnym modelu normalnym dostateczna i zupełna.
Statystyka $S^2=(n-1)^{-1}\sum_i(X_i-\bar X)^2$
ma rozkład niezależny od $\mu$; jest zatem poboczna dla
\emph{tego konkretnego modelu}. Twierdzenie Basu daje
niezależność $\bar X$ i $S^2$; w rozdziale 2 udowodniliśmy
ją już bezpośrednio przez obrót ortogonalny wektora normalnego.
Gdy $\sigma^2$ jest również nieznana, $S^2$ nie jest poboczna,
ponieważ jej rozkład zależy od $\sigma^2$. Nie wolno więc
mechanicznie przenosić argumentu między różnymi modelami.
\end{przyklad}


\begin{twierdzenie}[Maksimum próby jednostajnej jest zupełne]
\label{thm:uniform-complete}
Niech $n\ge2$ i $X_1,\ldots,X_n\iid U(0,\theta)$, $\theta>0$.
Wówczas $M=X_{(n)}$ jest statystyką dostateczną i zupełną.
Co więcej, statystyka $A=X_{(1)}/X_{(n)}$ jest poboczna;
na mocy twierdzenia Basu $A$ i $M$ są niezależne.
\end{twierdzenie}
\begin{proof}
Dla dodatnich obserwacji gęstość próby jest
$\theta^{-n}\ind{\{M\le\theta\}}$, więc faktoryzacja dowodzi
dostateczności $M$. Jego gęstość wynosi
$f_{M,\theta}(t)=n t^{n-1}\theta^{-n}\ind{\{0<t<\theta\}}$.
Niech $a$ spełnia $\E_\theta|a(M)|<\infty$ i
$\E_\theta a(M)=0$ dla każdego $\theta>0$.
Po pomnożeniu przez $\theta^n/n$ otrzymujemy
$G(\theta)=\int_0^\theta a(t)t^{n-1}\,\dd t=0$.
Funkcja $G$ jest lokalnie absolutnie ciągła na $(0,\infty)$,
gdyż $a(t)t^{n-1}$ jest lokalnie całkowalna; z twierdzenia
o różniczkowaniu całki mamy
$G'(\theta)=a(\theta)\theta^{n-1}=0$ dla niemal każdego
$\theta>0$. Zatem $a(M)=0$ $P_\theta$-p.n. dla każdego
$\theta$ i rodzina jest zupełna.
Ponadto przy zapisie $X_i=\theta U_i$, gdzie $U_i\iid U(0,1)$,
$A=U_{(1)}/U_{(n)}$ ma rozkład niezależny od $\theta$;
jest poboczna. Basu daje niezależność $A$ i $M$.
\end{proof}
\begin{uwaga}[Zaskakująca niezależność dwóch statystyk]
Maksimum i iloraz minimum do maksimum powstają z tych samych
obserwacji, a mimo to są niezależne w tym szczególnym modelu.
Twierdzenie nie mówi, że minimum i maksimum są niezależne
(na ogół nie są). Przy innych rozkładach lub niedokładnym
założeniu jednostajności omawiana niezależność może zniknąć.
\end{uwaga}

\section{Jaką metodę wybrać w konkretnym zadaniu?}
\begin{center}\small
\begin{tabular}{>{\raggedright\arraybackslash}p{0.25\textwidth}>{\raggedright\arraybackslash}p{0.31\textwidth}>{\raggedright\arraybackslash}p{0.31\textwidth}}
\toprule
\textbf{Problem} & \textbf{Najpierw sprawdź} & \textbf{Narzędzie i ograniczenie}\\
\midrule
Zredukować przechowywane dane & Czy istnieje statystyka dostateczna? & Faktoryzacja Fishera--Neymana; tylko w ustalonym modelu\\
Poprawić estymator nieobciążony & Czy znamy statystykę dostateczną? & Rao--Blackwell; poprawa MSE pod stratą kwadratową\\
Uzyskać UMVU & Czy statystyka jest \emph{zupełna i dostateczna}? & Lehmann--Scheffé; nie optymalizuje MSE wśród wszystkich obciążonych\\
Sprawdzić zbędność kolejności obserwacji & Czy próba rzeczywiście jest iid? & W modelach Bernoulliego i Poissona wystarcza suma; przy trendzie nie\\
Sprawdzić niezależność dwóch statystyk & Czy jedna jest zupełna dostateczna, druga poboczna? & Twierdzenie Basu; brak założeń wyklucza wniosek\\
Przewidzieć przyszłe zdarzenie & Jaki funkcjonał parametru chcemy oszacować? & UMVU lub estymator przez podstawienie MLE; zależnie od straty\\
\bottomrule
\end{tabular}
\end{center}

\paragraph{Granice stosowania.}
Dostateczność i zupełność są ściśle zależne od rodziny rozkładów.
Zachowanie sumy sukcesów jest bezstratne w modelu iid Bernoulliego,
lecz niekoniecznie przy zależności pomiędzy próbami lub niejednakowych
prawdopodobieństwach sukcesu. Ponadto \emph{najlepszy nieobciążony}
nie znaczy \emph{najlepszy przy dowolnej funkcji straty}.
Gdy wymagana jest kontrola ryzyka najgorszego przypadku,
wybór może paść na regułę minimaksową z poprzedniego rozdziału;
gdy najważniejsza jest odporność na odstające obserwacje,
na estymator odporny. W teorii statystycznej nie istnieje
optymalność niezależna od modelu i kryterium.

\Needspace{16\baselineskip}
\noindent\textbf{Przykład kodu: estymator UMVU dla $p^2$.}\par
\par
\Needspace{13\baselineskip}
\smallskip
\begin{lstlisting}[language=Python,caption={UMVU dla kwadratu prawdopodobieństwa sukcesu (Python).}]
x = [1, 0, 1, 0]
n, S = len(x), sum(x)
p2_umvu = S * (S - 1) / (n * (n - 1))
p2_plugin = (S/n)**2
print(p2_umvu, p2_plugin)  # 0.166666..., 0.25
\end{lstlisting}
\par
\Needspace{13\baselineskip}
\smallskip
\begin{lstlisting}[language=Julia,caption={UMVU dla kwadratu prawdopodobieństwa sukcesu (Julia).}]
x = [1, 0, 1, 0]
n, S = length(x), sum(x)
println((S*(S-1)/(n*(n-1)), (S/n)^2))
\end{lstlisting}

\paragraph{Podsumowanie.}
Metoda momentów i MLE \emph{konstruują} estymatory; teoria
dostateczności mówi, które aspekty danych zachowują informację
o parametrze; Rao--Blackwell poprawia nieobciążone estymatory,
a zupełność wraz z twierdzeniem Lehmanna--Scheffégo pozwala
uzasadnić ich optymalność w odpowiedniej klasie. Rozdział następny
poświęcimy informacji Fishera i dolnym ograniczeniom wariancji,
a następnie statystycznym twierdzeniom asymptotycznym.


% ================================================================
% ROZDZIAL 7 -- INFORMACJA FISHERA I DOLNE GRANICE
% ================================================================
